\documentclass[10pt,a4paper]{amsart}
\usepackage[margin=19mm]{geometry}
\usepackage{amsmath,amssymb,mathtools}
\usepackage[scr=rsfs,cal=euler]{mathalfa}
\usepackage{xfrac}
\usepackage[unicode,hidelinks,bookmarks=false]{hyperref}
\newcommand{\ie}{i.e.\ }
\newcommand{\cf}{cf.\ }
\newcommand{\resp}{respectively\ }
\newcommand{\Subsection}{\subsection}
\providecommand{\nobreakdash}{\nobreak\hbox{-}}
\setlength{\emergencystretch}{2em}
\multlinegap0pt
\usepackage{bm}
\usepackage{bbm}
\usepackage{dsfont}
\usepackage{xspace}
\usepackage{cancel}
\usepackage{mathtools}
\usepackage{cleveref}
\usepackage{amsthm}
\makeatletter
\@ifundefined{theorem}{\newtheorem{theorem}{Theorem}}{}
\@ifundefined{lemma}{\newtheorem{lemma}[theorem]{Lemma}}{}
\@ifundefined{proposition}{\newtheorem{proposition}[theorem]{Proposition}}{}
\makeatother
\newcommand{\adj}{^{\rm adj}}
\newcommand{\cof}{^{\rm cof}}
\newcommand{\mat}[1][n]{\operatorname{Mat}_{#1}(\mathbb{R})}
\newcommand{\trans}{^\top}
\title[The Strong Spectral Property and the Jacobian Method for Wei]{The Strong Spectral Property and the Jacobian Method for Weighted Laplacian Matrices}
\author{Minerva Catral, Shaun Fallat, Himanshu Gupta, Jephian C.-H. Lin}
\date{Source: arXiv:2602.18999v1 (CC BY 4.0)}
\begin{document}
\maketitle

\noindent\emph{Source and licence.} The Strong Spectral Property and the Jacobian Method for Weighted Laplacian Matrices. Version arXiv:2602.18999v1 (CC BY 4.0), distributed under the Creative Commons Attribution 4.0 International licence, as recorded in the arXiv licence metadata for this exact version and shown on the abstract page https://arxiv.org/abs/2602.18999v1. Full rights notice: licenses/arxiv-2602.18999-CC-BY-4.0.txt.\par\smallskip

\noindent\textit{Editorial scope.} Counted unit: the Lemma whose source label is \texttt{lem:gradientind} in the article (source0.tex lines 637--642, source label \texttt{lem:gradientind}), delivered together with the article's own complete original proof of it (source0.tex lines 643--679, one \texttt{proof} environment of 37 lines). No mathematics is written, repaired or re-derived: statement and proof are the article's own text line for line. Required context retained uncounted: prop:gradientadj (lines 599--605). Every definition, display and label of those lines is retained unchanged. Omitted from this package and not redistributed: 1--598, 606--636, 680--1585. Nothing inside a retained excerpt is omitted or altered: every retained body is delivered exactly as the article's lines are written, so no omission receipt of any kind accompanies this package. Citations: the retained excerpts carry no citation command at all, so no bibliography entry of the article is transcribed and no reference is redistributed. Reference adaptation: none was needed and none was made; every reference inside the delivered unit resolves inside it, so no unresolved reference or citation is produced. Printed slips kept literal: no printed word, symbol, hypothesis or number of the counted statement or of its proof was altered. Layout and class: the article's own class is \texttt{article} with packages including inputenc, fontenc, amsmath, amssymb, amsthm, ulem, hyperref, cite, cleveref, enumitem, xcolor, tikz, soul, cancel; the wrapper is the lane's pinned \texttt{amsart} editorial wrapper at 10pt on A4, which restates the theorem-like environments the retained text needs and adds the article's own macro definitions that the retained text needs (\texttt{\textbackslash adj}, \texttt{\textbackslash cof}, \texttt{\textbackslash mat}, \texttt{\textbackslash trans}), with extra wrapper packages (bm, bbm, dsfont, xspace, cancel, mathtools, cleveref). The delivered numbering is the unit's own. Assets: no retained span contains a figure environment, a graphics-inclusion command, a PDF-page inclusion, a table or a TikZ picture; no image file is copied into this package, the package declares no asset dependency and no image file is redistributed with it. Licence: version arXiv:2602.18999v1 of this article is distributed under the Creative Commons Attribution 4.0 International licence, as recorded in the arXiv licence metadata for this exact version and shown on the abstract page https://arxiv.org/abs/2602.18999v1. A full rights notice is delivered at licenses/arxiv-2602.18999-CC-BY-4.0.txt. Characters: every retained excerpt is pure ASCII, so the delivered engine, XeLaTeX with its default Latin Modern text font, reports no missing character.
\begin{proposition}
\label{prop:gradientadj}
Let $M\in\mat$ and $\nabla s_k = \nabla s_k(M)$ for $k = 1, \ldots, n$.  Then the total derivative of $\det(M - xI)$ with respect to the entries is 
\[
    \nabla\det(M - xI) = \nabla s_1(-x)^{n-1} + \nabla s_2(-x)^{n-2} + \cdots + \nabla s_n = (M - xI)\cof.
\] 
\end{proposition}

\begin{lemma}
\label{lem:gradientind}
Let $M\in\mat$ and $\nabla s_k = \nabla s_k(M)$ for $k = 1, \ldots, n$.  Then $
\{\nabla s_1, \ldots, \nabla s_n\}$ and $\{I, M\trans, \ldots, (M\trans)^{n-1}\}$ span the same subspace.  Consequently, $
\{\nabla s_1, \ldots, \nabla s_n\}$ is linearly independent if and only if the minimal polynomial of $M$ has degree $n$.  
\end{lemma}

\begin{proof}
To simplify notation, let $M_x = M - xI$, treated as a matrix over the field of rational functions $\mathbb{R}(x)$, and consider its characteristic polynomial 
\[
    \det(M_x - tI) = p_0(x)(-t)^n + p_1(x)(-t)^{n-1} + \cdots + p_n(x).
\]
By the Cayley--Hamilton theorem, the inverse of $M_x$ can be expressed as 
\[
    M_x^{-1} = \frac{1}{p_n(x)}(p_0(x)(-M_x)^{n-1} + p_1(x)(-M_x)^{n-2} + \cdots + p_{n-1}(x)I).
\]
Now, it is known that 
\[
    M_x\cof = (M_x\adj)\trans = \det(M_x)(M_x\trans)^{-1}.
\]
Thus, since $p_0(x) = 1$ and $p_n(x) = \det(M_x)$, we have 
\[
    M_x\cof = (-M_x\trans)^{n-1} + p_1(x)(-M_x\trans)^{n-2} + \cdots + p_{n-1}(x)I.
\]
Expanding the powers of $-M_x\trans$ on the right-hand side, we may write 
\[
    M_x\cof = r_0(x)(M\trans)^{n-1} + r_1(x)(M\trans)^{n-2} + \cdots + r_{n-1}(x)I
\]
for some coefficients $r_0(x), \dots, r_{n-1}(x)$.
Equating this expression  for $M_x\cof$ with the formula in \cref{prop:gradientadj}, each element in $\{\nabla s_1, \ldots, \nabla s_n\}$ is in the span of $\{I, M\trans, \ldots, (M\trans)^{n-1}\}$.  

On the other hand, we claim that $r_k(x)$ is a polynomial of degree $k$.  This would imply that $\{r_0(x), \ldots, r_{n-1}(x)\}$ is a basis of the space of polynomials of degree less than $n$, and  therefore, each element of $\{1, (-x), \ldots, (-x)^{n-1}\}$ can be written as a linear combination of $\{r_0(x), \ldots, r_{n-1}(x)\}$.  Consequently, by comparing the two expressions for  $M_x\cof $, each element in $\{I, M\trans, \ldots, (M\trans)^{n-1}\}$ is  in the span of $\{\nabla s_1, \ldots, \nabla s_n\}$, showing that the  two sets  have the same span.   We now examine the leading term of $r_k(x)$.    By direct computation, 
\[
    r_k(x) = \sum_{i=0}^{k} p_i(x) \cdot (-1)^{n-1-i} \binom{n-1-i}{k-i}(-x)^{k-i}.
\]
Note that $p_i(x)$ is the sum of all $i\times i$ principal minors of $M_x$, so its leading term is $\binom{n}{i}(-x)^i$.  Thus, the leading term of $r_k(x)$ is 
\[
    (-x)^{k}\sum_{i=0}^{k} \binom{n}{i} \cdot (-1)^{n-1-i} \binom{n-1-i}{k-i} = (-1)^{n-1}x^{k}\sum_{i=0}^{k} \binom{n}{i}\binom{-(n-k)}{k-i}.
\]
Observe that the summation is the coefficient of $y^k$ in $(1+y)^n(1+y)^{-(n-k)} = (1+y)^k$, which is $1$.   Therefore, the leading term of $r_k(x)$ is $(-1)^{n-1}x^k$.  This completes our argument showing that $
\{\nabla s_1, \ldots, \nabla s_n\}$ and $\{I, M\trans, \ldots, (M\trans)^{n-1}\}$ span the same subspace.

If the minimal polynomial is of degree less than $n$, then the dimension of the span of $\{I, M\trans, \ldots, (M\trans)^{n-1}\}$ is less than $n$, indicating that $\{\nabla s_1, \ldots, \nabla s_n\}$ is linearly dependent.  If  the minimal polynomial is of degree $n$, then the dimension of the span of $\{I, M\trans, \ldots, (M\trans)^{n-1}\}$ is $n$, implying that $\{\nabla s_1, \ldots, \nabla s_n\}$ is a linearly independent set.
\end{proof}

\end{document}
